\section{Task 1A: Understanding the IP-Stride Prefetcher}
\label{sec:task1a}
\subsection{State and lookup}
The supplied \texttt{prefetcher/}\allowbreak\texttt{ip\_stride.l1d\_pref} implementation defines \texttt{IP\_TRACKER\_COUNT=24}, \texttt{PREFETCH\_DEGREE=3}, and \texttt{MASK\_IP=0xFF}. Each tracker contains \texttt{ip}, \texttt{last\_cl\_addr}, \texttt{last\_stride}, and an LRU rank. The IP is masked to its lower eight bits before lookup, so instruction pointers that differ only above those bits alias to the same logical IP. The implementation does not use a valid bit in its tracker structure.

On an IP miss, it chooses the tracker whose LRU rank is 23, stores the masked IP and current cache-line address, resets the previous stride to zero, updates LRU ranks, and returns without prefetching. On a hit, it calculates a signed stride in cache-line units as current block address minus the previous block address. A zero stride returns immediately. Otherwise, prefetches are issued only when the current stride equals the previously recorded stride. This is a two-consecutive-equal-strides test rather than a saturating confidence counter.

\subsection{Prefetch generation and safety}
The next three block addresses are computed as $A + k\Delta$, for $k=1,2,3$, and converted to byte addresses by shifting by \texttt{LOG2\_BLOCK\_SIZE}. A candidate is discarded, and the loop stops, if it is outside the same 4 KiB page as the current demand access. If MSHR occupancy is below half of MSHR capacity, the prefetch is requested to fill L1; otherwise it is requested to fill L2. The final-statistics callback prints the fixed degree but does not report useful/issued totals.

\subsection{Example state evolution}
The following example uses cache-block numbers (not byte addresses) and assumes every address remains in the same page.
\begin{table}[H]\centering\small
\caption{Example for one masked IP and block sequence 100, 102, 104, 106.}
\begin{tabularx}{\linewidth}{c c c X X}
\toprule Access & Block & Observed stride & Stored state after access & Prefetch action \\
\midrule
1 & 100 & -- & Last block = 100; last stride = 0 & Allocate tracker; none \\
2 & 102 & $+2$ & Last block = 102; last stride = $+2$ & No: previous stride was 0 \\
3 & 104 & $+2$ & Last block = 104; last stride = $+2$ & Yes: blocks 106, 108, 110 \\
4 & 106 & $+2$ & Last block = 106; last stride = $+2$ & Yes: blocks 108, 110, 112 \\
\bottomrule
\end{tabularx}\end{table}

\subsection{Strengths and limitations}
The design is simple and inexpensive, and it can issue multiple requests as soon as a constant stride repeats. It does not encode a multi-stride signature, however, so a repeating sequence such as $+1,+2,+2,+2,+3$ is not modeled as a cycle. It may fail to prefetch across a page boundary even if the stream continues contiguously, and it can lose history when a tracker is evicted. Because tracker IPs are masked and there is no valid bit, a masked IP equal to zero can also match an initially zero-initialized tracker. These limitations motivate the extra history and signature state in Task 1B.
